\subsection{Lexikální analýza}

Lexikální analyzátor (\texttt{scanner.c}) je deterministický konečný automat.
Aktuální stav drží proměnná typu \texttt{State}; každému nekoncovému uzlu
diagramu v~příloze~\ref{app:fsm} odpovídá jedna hodnota tohoto výčtu
a v~komentáři u~ní je uveden název uzlu z~diagramu. Parser volá funkci
\texttt{get\_token()}, která čte vstup po znacích, prochází automatem a při
každém volání vrátí jeden token. Text lexému se během čtení ukládá do
dynamického řetězce (\texttt{dynstr}), takže délka identifikátorů ani
řetězců není omezena.

\paragraph{Rozhraní s~parserem.}
Token (\texttt{Token}) obsahuje typ, atribut (řetězec nebo \texttt{NULL}),
číslo řádku a příznak \texttt{newline\_before} (viz níže). Atribut mají jen
identifikátory a literály:
\begin{itemize}

  \item číselný literál zůstává jako text, jen bez podtržítek a bez přípony
        \texttt{d}/\texttt{D}; na číslo ho převádějí až další části překladače,
  \item u~řetězcového literálu jsou escape sekvence už nahrazeny skutečnými
        znaky,
  \item klíčové slovo se pozná až po načtení celého identifikátoru, a to
        vyhledáním v~tabulce klíčových slov; rezervována jsou i klíčová slova
        rozšíření.
\end{itemize}
Například řádek \texttt{val x = 3.5} skener rozloží na tokeny
\texttt{T\_VAL}, \texttt{T\_ID}\,(\texttt{x}), \texttt{T\_ASSIGN},
\texttt{T\_LIT\_DOUBLE}\,(\texttt{3.5}) a \texttt{T\_EOL};
v~závorce je uveden atribut.

Samotné podtržítko, samotné \texttt{|} a hranaté závorky \texttt{[},
\texttt{]} jsou ve specifikaci určeny pro rozšíření. Skener je vrací jako
tokeny (\texttt{T\_UNDERSCORE}, \texttt{T\_PIPE}, \texttt{T\_LBRACKET},
\texttt{T\_RBRACKET}) a jejich nepovolené použití odhalí parser jako
syntaktickou chybu. Paměť atributu uvolňuje ten, kdo token převzal, funkcí
\texttt{free\_token()}. Po posledním tokenu (i po chybě) se volá
\texttt{scanner\_free()}, která uvolní zásobník odsazení.

\paragraph{Odsazení a konce řádků.}
Na konci každého neprázdného řádku vrátí skener token \texttt{T\_EOL}.
Na začátku dalšího řádku spočítá mezery a porovná je s~úrovní aktuálního
bloku (vrchol zásobníku odsazení):
\begin{itemize}
  \item stejné odsazení -- pokračuje se ve stejném bloku,
  \item větší odsazení -- začíná nový blok, skener vrátí \texttt{T\_INDENT}
        a novou úroveň uloží na zásobník,
  \item menší odsazení -- končí jeden nebo více bloků, skener vrací
        \texttt{T\_DEDENT} (při každém volání jeden), dokud se nedostane na
        stejnou úroveň. Pokud taková úroveň na zásobníku není (odsazení
        padne mezi dvě uložené úrovně), jde o~lexikální chybu.
\end{itemize}
Tabulátor v~odsazení je lexikální chyba. Prázdné řádky a řádky jen
s~komentářem odsazení neovlivňují. Na konci souboru skener vrátí
\texttt{T\_EOL} (i když poslední řádek nekončí znakem konce řádku),
potom \texttt{T\_DEDENT} pro každý ještě otevřený blok a nakonec
\texttt{T\_EOF}.

Uvnitř kulatých závorek se konec řádku bere jako mezera, takže výraz nebo
seznam argumentů může pokračovat na dalším řádku. Takový řádek musí být
odsazen víc než začátek příkazu; jen uzavírací \texttt{)} smí být na stejné
úrovni. Pokud byl uvnitř závorek před tokenem přeskočen konec řádku, má
token nastaven příznak \texttt{newline\_before}. Podle něj parser pozná,
zda je povolena čárka před \texttt{)}.

\paragraph{Odchylky od konečného automatu.}
Kromě aktuálního stavu si skener pamatuje ještě několik údajů, které čistý
konečný automat nemá:
\begin{itemize}
  \item \textbf{zásobník odsazení} -- pamatuje si úrovně odsazení
        otevřených bloků, aby skener poznal, kdy vrátit \texttt{T\_INDENT}
        a kolik \texttt{T\_DEDENT}; podle potřeby se zvětšuje,
  \item \textbf{hloubku vnoření komentářů} $h$ -- blokové komentáře lze
        vnořovat, a komentář tak končí až u~posledního \texttt{*/}.
        Zároveň si pamatuje, zda komentář obsahoval konec řádku; takový
        komentář se pak chová jako konec řádku,
  \item \textbf{počet otevřených kulatých závorek} a to, zda byl uvnitř
        nich přeskočen konec řádku (viz výše),
  \item \textbf{jeden znak navíc k~vrácení} -- u~zápisu \texttt{3.toString}
        se skener dozví, že nejde o~desetinné číslo, až po přečtení
        \texttt{3}, tečky a dalšího znaku. Vrátí tedy \texttt{3} jako
        \texttt{Int} a tečku i další znak vrátí na vstup; protože
        \texttt{ungetc} zaručuje vrácení jen jednoho znaku, druhý si skener
        drží ve vlastní proměnné,
  \item \textbf{klíčová slova} se v~automatu nerozlišují od identifikátorů;
        po načtení celého slova se vyhledá v~tabulce klíčových slov,
  \item \textbf{slovo za \texttt{@}} se také načte celé a porovná
        s~\texttt{main}, místo aby mělo vlastní řetězec stavů,
    \item \textbf{rozsah desetinného literálu} se kontroluje funkcí
        \texttt{strtod} až po načtení celého čísla; příliš velké číslo
        i číslo, které by se zaokrouhlilo na nulu, je lexikální chyba
        (nad rámec zadání, specifikace tuto kontrolu nevyžaduje).
\end{itemize}

\paragraph{Chyby.}
Lexikální chybou je například neznámý znak, chybný zápis čísla
(\texttt{05} -- počáteční nuly, \texttt{1\_}, \texttt{1e}), neukončený
řetězec nebo komentář, neplatná escape sekvence, desetinný literál mimo
rozsah nebo chybné odsazení.
Skener v~takovém případě vypíše na \texttt{stderr} hlášení s~číslem řádku
a \texttt{get\_token()} vrátí~1. Parser tento kód předá dál a překladač
skončí s~návratovým kódem~1. Pokud selže alokace paměti, překladač skončí
s~kódem~99.